\subsection{Method: From Numerical Methods to AI \quad / \quad \textthai{วิธีการ: จากวิธีการเชิงตัวเลขสู่ AI}}
\begin{paracol}{2}
Traditionally, we solve these equations using numerical integration methods like the \textbf{Euler Method} or \textbf{Runge-Kutta (RK4)}. While efficient for single oscillators, these methods scale poorly when simulating billions of interacting particles or high-dimensional quantum systems.

AI changes this by learning the \textbf{evolution mapping}\index{Evolution Mapping@แผนภูมิวิวัฒนาการ (Evolution Mapping)} directly. Instead of integrating step-by-step, we train a neural network to predict the state at any future time $t$, effectively "learning" the continuous solution $x(t)$.

\switchcolumn

\begin{thai}
ตามประเพณีดั้งเดิม เราแก้สมการเหล่านี้โดยใช้วิธีการรวมกลุ่มเชิงตัวเลข เช่น \textbf{Euler Method} หรือ \textbf{Runge-Kutta (RK4)} แม้ว่าจะมีประสิทธิภาพสำหรับตัวสั่นพ้องตัวเดียว แต่วิธีการเหล่านี้ปรับขนาดได้ยากเมื่อต้องจำลองอนุภาคนับพันล้านที่มีปฏิสัมพันธ์กัน หรือระบบควอนตัมที่มีมิติสูง

AI เปลี่ยนแปลงสิ่งนี้โดยการเรียนรู้ \textbf{Evolution Mapping} โดยตรง แทนที่จะรวมกลุ่มทีละขั้นตอน เราฝึกโครงข่ายประสาทเทียมเพื่อทำนายสถานะ ณ เวลาใดๆ ในอนาคต $t$ ซึ่งเป็นการ "เรียนรู้" คำตอบต่อเนื่อง $x(t)$ ได้อย่างมีประสิทธิภาพ
\end{thai}
\end{paracol}
